\par\medskip
\begin{center}\bfseries
V. Die nachgelassene Abhandlung über das agM.\\
Ausbau zur Theorie der Modulfunktion.
\end{center}
\begin{center}\textls[100]{a) Einleitendes.}\end{center}
\noindent\hspace*{\parindent}Dass \textsc{Gauss} in der Tagebuchnotiz 108. von der »Entwicklung aller nur
denkbaren lemniskatischen Funktionen« spricht, zeigt, dass er sich bewusst
war, durch seine Untersuchungen zwei verschiedene Probleme der Lösung zu-\par
\clearpage\SourceLeaf{170}\ArithmeticHead{85}{R}{Va. EINLEITENDES.}
\noindent geführt zu haben. Erstens hat er, ausgehend von einer Grösse \(x\), die absolut
genommen kleiner ist als Eins, die Funktion
\[
\operatorname{Sl}\varphi=\frac{q(x)}{r(x)}\cdot\frac{\vartheta_{11}(\psi|x)}{\vartheta_{00}(\psi|x)},\quad
\psi=\frac{\varphi}{\pi q(x)^2},
\]
\noindent hergestellt, die die Umkehrung des Integrals
\begingroup\makeatletter\tagsleft@true\let\veqno\leqno\makeatother
\begin{equation*}
\varphi=\int_0^S\frac{dS}{\sqrt{(1-S^2)(1+\mu^2S^2)}}\tag{1}
\end{equation*}
\endgroup
\noindent für \(\mu=r(x)^2/q(x)^2\) liefert\footnote{Neben der Funktion Sl bildet \textsc{Gauss} in der Scheda Ac, Werke X, 1, S. 200 noch, wie für die spezielle lemniskatische Funktion \(\mu=1\) (vergl. oben S. 33),
\[
\operatorname{cl}\varphi=\operatorname{Sl}(\tfrac12\varpi-\varphi)=\frac qr\frac{\vartheta_{10}(\psi|x)}{\vartheta_{01}(\psi|x)}.
\]
Nach der Analogie der \textsc{Jacobischen} Funktionen
\[
\begin{gathered}
\operatorname{sn}\varphi=\frac pr\frac{\vartheta_{11}(\chi|x)}{\vartheta_{01}(\chi|x)},\quad
\operatorname{cn}\varphi=\frac qr\frac{\vartheta_{10}(\chi|x)}{\vartheta_{01}(\chi|x)},\\
\operatorname{dn}\varphi=\frac qp\frac{\vartheta_{00}(\chi|x)}{\vartheta_{01}(\chi|x)},\quad
\chi=\frac{\varphi}{\pi p(x)^2},
\end{gathered}
\]
wollen wir setzen
\[
\begin{gathered}
\operatorname{Cl}\varphi=(1-\operatorname{Sl}^2\varphi)^{\frac12}=\frac pr\frac{\vartheta_{10}(\psi|x)}{\vartheta_{00}(\psi|x)},\\
\operatorname{Dl}\varphi=(1+\mu^2\operatorname{Sl}^2\varphi)^{\frac12}=\frac pq\frac{\vartheta_{01}(\psi|x)}{\vartheta_{00}(\psi|x)};
\end{gathered}
\]
diese Funktionen finden sich tatsächlich bei \textsc{Gauss}, Werke III, S. 473 in der Aufzeichnung von 1827, Art. 5.}; zweitens aber war für ein gegebenes \(\mu^2\) die Umkehrungsfunktion des Integrals (1) zu bestimmen. Wir haben im vorhergehenden Abschnitt unser Augenmerk ausschliesslich auf das erste Problem
gerichtet und sind darum auch in der Form der Faktoren, die in den Formeln
für die Transformation der Thetafunktionen auftreten, von \textsc{Gauss'} Darstellung
in der Scheda Ac abgewichen, was übrigens \textsc{Gauss} selbst in seinen spätern
Darstellungen (z. B. Werke III, S. 472, vom Jahre 1827) getan hat. Dabei
haben wir aber schon auf die Frage hingewiesen, die zu lösen bleibt, wenn
man das zweite Problem in Angriff nehmen will. Die Aufgabe, zu einem
gegebenen Integral (1), d. h. zu einem gegebenen Werte von \(\mu=\operatorname{tg}v\), die
Perioden \(\varpi,\varpi'\)\footnote{In Wirklichkeit sind \(2\varpi\) und \(2i\varpi'\) die Perioden.} anzugeben, führt zunächst auf die Bestimmung der beiden
agM. \(M(1,\cos v),M(1,\sin v)\). Es handelt sich dabei wesentlich darum, zur
Berechnung dieser Mittel auch für beliebige komplexe Werte von \(\mu\) brauchbare Methoden anzugeben. Wir kommen darauf im Anschluss an den Bericht über die schon öfter erwähnte Abhandlung aus dem Jahre 1800\footnote{\textit{De origine proprietatibusque generalibus numerorum arithmetico-geometricorum}, Handbuch 15, Ba,
abgedruckt Werke III, S. 361—374, deutsch bei \textsc{Geppert}, a. a. O.\textsuperscript{83)}, S. 27ff. \textsc{Gauss} schreibt an \textsc{Schumacher} im April 1816 (Werke X, 1, S. 247): »Ich habe zwar \textls[100]{ausser} jenem [nämlich dem richtigen agM.] auch noch andere arithmetisch-geometrische Mittel betrachtet, die aber ganz elementarisch sind. Jenes ist das wahre, worüber Sie hier [nämlich in Göttingen, wo sich \textsc{Schumacher} 1808/09 studienhalber aufgehalten hat] auch eine im Jahre 1800 von mir angefangene kleine Abhandlung gelesen haben, in einem blauen Oktavbande, Varia betitelt, worin noch von \textls[100]{Ihrer} Hand eine Restitutio in Integrum einiger durch einen Tintenfleck unkenntlich gewordener Stellen zu sehen ist.« Dieser »blaue Oktavband« ist das jetzt mit 15, Ba, bezeichnete Handbuch, die 1800 begonnene Abhandlung, die im Text genannte, in deren Handschrift die Spuren der von \textsc{Gauss} erwähnten Tintenflecke noch deutlich zu erkennen sind.
Im Artikel 7 dieser Abhandlung (Werke III, S. 368) wird der Artikel 162 der \textit{Disquisitiones arithmeticae} (erschienen 1801) zitiert. Man braucht darum der von \textsc{Gauss} gegebenen Datierung kein Misstrauen entgegenzubringen; denn, wie \textsc{Gauss} an \textsc{Bolyai} schreibt (siehe \textit{Briefwechsel Gauss-Bolyai}, S. 35), hat er am 16. Dezember 1799 in Helmstedt die Korrektur des 18. Bogens der \textit{Disqu. arithm.} erhalten, der Artikel 162 (S. 176 der Originalausgabe) war also 1800 längst gedruckt. Ein Fragment einer deutschen Redaktion der Abhandlung über das agM. findet sich im Nachlass (Ff), es enthält einen Teil des Art. 2 (Werke III, S. 362, von Zeile 16 an), dann einige Beispiele (nämlich \(a=1,b=0{,}2\), wie im Art. 3, aber mit weniger Dezimalstellen, und \(a=5,b=2\)), ferner Art. 4 und endlich flüchtig angedeutet den Art. 9 des zweiten Teils, abschliessend mit der Formel S. 373 unten: \(dM(x,y)=\mathrm{etc.}\)}:\par
\clearpage\SourceLeaf{171}\ArithmeticHead{86}{L}{L. SCHLESINGER, ÜBER GAUSS’ FUNKTIONENTHEORETISCHE ARBEITEN.}
\noindent \textit{Über den Ursprung und die allgemeinen Eigenschaften der arithmetisch-geometrischen Mittelzahlen} weiter unten zurück.

Mit der blossen Herstellung der \(\varpi,\varpi'\) ist jedoch das in Rede stehende
Problem noch nicht erledigt. Es handelt sich vielmehr, wie schon oben hervorgehoben wurde, noch um die Bildung der Thetafunktionen \(T\) und \(W\), die
erfordert, dass zur Sicherung der Konvergenz der in Betracht kommenden
Reihen der reelle Teil von
\begingroup\makeatletter\tagsleft@true\let\veqno\leqno\makeatother
\begin{equation*}
t=\frac{\varpi'}{\varpi}=\frac{M(1,\cos v)}{M(1,\sin v)}=-\frac{\log x}{\pi}\tag{2}
\end{equation*}
\endgroup
\noindent einen positiven Wert habe. \textsc{Gauss} war also in jenen ersten Junitagen des
Jahres 1800 jedenfalls bestrebt, der hiermit bezeichneten Schwierigkeit Herr
zu werden. Da die Frage wesentlich der Theorie des agM. angehört, wird
er vor allem bemüht gewesen sein, diese Theorie weiter zu vertiefen, und zwar
nach der bis dahin noch nicht erörterten Seite hin, dass durch die Willkürlichkeit des bei der Bildung der geometrischen Mittel auftretenden Vorzeichens
eine Mehrdeutigkeit des Algorithmus bedingt wird. Setzt man
\begingroup\makeatletter\tagsleft@true\let\veqno\leqno\makeatother
\begin{equation*}
\sin v=k,\quad \cos v=(1-k^2)^{\frac12}=k',\tag{3}
\end{equation*}
\endgroup
\noindent so wird die Frage präziser so zu stellen sein: gibt es für ein beliebig vorgeschriebenes \(k\) Determinationen der beiden agM. \(M(1,k),M(1,k')\), für die\par
\clearpage\SourceLeaf{172}\ArithmeticHead{87}{R}{Va. EINLEITENDES.}
\noindent der reelle Teil von
\begingroup\makeatletter\tagsleft@true\let\veqno\leqno\makeatother
\begin{equation*}
t=\frac{M(1,k')}{M(1,k)}\tag{2\,a}
\end{equation*}
\endgroup
\noindent positiv ausfällt, und wie lassen sich diese »zusammengehörigen« Determinationen charakterisieren? Dass \textsc{Gauss} in der Tat diesen Weg eingeschlagen hat,
zeigt die Tagebuchnotiz 109., vom 3. Juni 1800 datiert, also in unmittelbarem
Anschluss an die in 108. geschilderten Ergebnisse geschrieben\footnote{Man vergleiche auch im Faksimile des \textit{Tagebuchs}, Werke X, 1, die Gleichheit der Schriftzüge in den Notizen 108. und 109.}, die so lautet:
»Zwischen zwei gegebenen Zahlen gibt es sowohl für das arithmetisch-geometrische als für das harmonisch-geometrische Mittel\footnote{Das harmonische Mittel zwischen zwei Zahlen \(a,b\) ist die durch die Gleichung \(a_1=\dfrac{2ab}{a+b}\) definierte Grösse. Demnach wäre der Algorithmus des harmonisch-geometrischen Mittels durch die Gleichungen
\[
\begin{gathered}
\frac1{a_{n+1}}=\frac12\left(\frac1{a_n}+\frac1{b_n}\right),\quad
\frac1{b_{n+1}}=\sqrt{\frac1{a_n}\cdot\frac1{b_n}}\\
(n=0,1,2,\ldots;\ a_0=a,\ b_0=b)
\end{gathered}
\]
definiert. Es ist also
\[
\lim_n\frac1{a_n}=\lim_n\frac1{b_n}=M\left(\frac1a,\ \frac1b\right)
\]
und folglich
\[
\lim_n a_n=\lim_n b_n=\frac{ab}{M(a,b)},
\]
wodurch das harmonisch-geometrische Mittel auf das agM. zurückgeführt ist (vergl. \textsc{Th. Lohnstein}, Zeitschrift für Math. u. Physik \textbf{30}, 1888, S. 316). In bezug auf dieses Mittel vergl. die Notiz aus Scheda Ac, Werke X, 1, S. 551.} unendlich viele Bestimmungen; dass wir deren gegenseitigen Zusammenhang in vollkommener
Weise durchschauen konnten, hat uns ein Glücksgefühl bereitet.«

\textsc{Gauss} will also damit sagen, dass er die unendlich vielen Werte von \(t\)
anzugeben imstande ist, die zu einem und demselben Wertepaar \(k,k'\) gehören,
d. h. in der uns geläufigen Ausdrucksweise, dass er für die \textls[100]{Modulfunktion}
\(k'\) als Funktion von \(t\), wie sie durch die Gleichungen (3), (2\,a) definiert wird,
die Fundamentaleigenschaft, nämlich ihren \textls[100]{Automorphismus} entdeckt habe.
Er wird demnach die Tatsache, dass für ein willkürliches \(\mu\) oder \(k'\) der reelle
Teil von \(t\) positiv ausfällt, wenn man die Determinationen der agM. geeignet
wählt, mit Hilfe der Theorie der Modulfunktion bewiesen haben. Und daraus
kann man wieder rückwärts auf die Tragweite der Notiz 109. schliessen. \textsc{Gauss}
muss in jenen Tagen erkannt haben, dass die inverse Funktion der Modul-\par
\clearpage\SourceLeaf{173}\ArithmeticHead{88}{L}{L. SCHLESINGER, ÜBER GAUSS’ FUNKTIONENTHEORETISCHE ARBEITEN.}
\noindent funktion nur Werte mit positivem Realteil annimmt, er muss sich also alle
wesentlichen Eigenschaften dieser Funktion entwickelt haben. Dass es sich am
3. Juni um eine grosse Entdeckung handelt, geht schon aus der Wendung
»hat uns ein Glücksgefühl bereitet« hervor; \textsc{Gauss} gebraucht diese oder eine
ähnliche Wendung nur sehr selten, sonst im \textit{Tagebuch} nur bei zahlentheoretischen Notizen (siehe die Nrn. 30, 73, 114, 141) und man wird darum vermuten können, dass er damals gleich auch den Zusammenhang mit der Transformationstheorie der binären quadratischen Formen durchschaut hat\footnote{In einem Brief an \textsc{Gerling} vom 6. Januar 1819 (Werke X, 1, S. 122) schreibt \textsc{Gauss}: »für mich wenigstens sind und bleiben die Untersuchungen der höheren Arithmetik bei weitem das Allerschönste der Mathematik, und der Genuss, den ich auch an den schönsten astronomischen Untersuchungen finde, ist gar Nichts, verglichen mit dem, welchen die höhere Arithmetik gewährt«.}. Wenn
natürlich auch die Möglichkeit, ja sogar die Wahrscheinlichkeit vorliegt,
dass die in diesen Tagen des Jahres 1800 gefundenen Resultate später ergänzt und vertieft worden sind, so dürfen wir doch die \textls[100]{Prinzipien} alles
dessen, was uns von \textsc{Gauss} über die Theorie der Modulfunktion erhalten ist,
für die hier in Rede stehende Zeit in Anspruch nehmen. Wir wollen also
diese Entwicklungen — auf die Gefahr hin, einzelne erst nach 1800 gefundene
Details schon hier einzuordnen — im Zusammenhange besprechen, zumal,
wie schon am Schluss des Abschnitts IV bemerkt worden ist, Aufzeichnungen
aus der Zeit der Tagebuchnotiz selbst fehlen.

Wir wenden uns nun zur Besprechung der mit der Tagebuchnotiz 109.
im Zusammenhang stehenden Aufzeichnungen und beginnen, ihres elementaren
Ausgangspunktes wegen, mit der nachgelassenen Abhandlung vom Jahre 1800.
